% Part: sets-functions-relations
% Chapter: size-of-sets-complete

\documentclass[../../../include/open-logic-chapter]{subfiles}

\begin{document}

\olchapter{sfr}{siz}{संचांचे आकारमान}

\begin{editorial}
या प्रकरणात प्रगणने, गणनीयता आणि अगणनीयता यांची चर्चा आहे.
काही विभागांच्या दोन आवृत्त्या आहेत: अधिक प्राथमिक आवृत्तीत प्रगणने
म्हणजे याद्या किंवा $\PosInt$ पासूनची आच्छादने मानली आहेत;
अधिक अमूर्त आवृत्तीत प्रगणनांची व्याख्या $\Nat$ बरोबरच्या
एकास-एक आच्छादनांनी केली आहे.
\end{editorial}

\olimport{introduction}

\olimport{enumerability}

\olimport{zig-zag}

\olimport{pairing}

\olimport{pairing-alt}

\olimport{non-enumerability}

\olimport{reduction}

\olimport{equinumerous-sets}

\olimport{comparing-size}

\olimport{schroder-bernstein}

\begin{editorial}
पुढील \olref[sfr][siz][enm-alt]{sec},
\olref[sfr][siz][nen-alt]{sec} आणि \olref[sfr][siz][red-alt]{sec}
हे अनुक्रमे \olref[sfr][siz][enm]{sec},
\olref[sfr][siz][nen]{sec} आणि \olref[sfr][siz][red]{sec}
यांचे पर्यायी विभाग आहेत. टिम बटन यांनी त्यांच्या Open Set Theory
या ग्रंथात वापरण्यासाठी ते तयार केले. ते किंचित अधिक प्रगत आहेत आणि
संच सिद्धांताच्या संदर्भात अधिक योग्य अशी गणनीयतेची भिन्न व्याख्या
वापरतात: यादी देता येणे किंवा $\PosInt$ पासूनच्या आच्छादक फलनाची
व्याप्ती असणे याऐवजी $\Nat$ किंवा त्याच्या आरंभीच्या खंडाशी एकास-एक
आच्छादन असणे.
\end{editorial}

\olimport{enumerability-alt}
\olimport{non-enumerability-alt}
\olimport{reduction-alt}

\OLEndChapterHook

\end{document}
